% 题干转写：完成
% 解析转写：完成
% 图形核对：完成
\section{第37届复赛第7题：LIGO双黑洞引力波辐射估算}
\subsection{题干}
在双黑洞并合的旋近阶段，将质量为 $M$、$m$ 的双黑洞近似为绕共同质心运动的质点，间距为 $L$，如图\ref{fig:cpho37q07}。系统仅通过引力波辐射损失能量，采用牛顿引力理论近似。
\begin{figure}[htbp]\centering
\begin{tikzpicture}[x=1cm,y=1cm,>=Stealth,font=\small]
  \coordinate (M) at (-2.35,0); \coordinate (O) at (0,0); \coordinate (m) at (2.35,0);
  \draw[thick] (M) -- (m);
  \fill (M) circle (2.2pt) node[above left] {$M$};
  \fill (O) circle (1.5pt) node[above] {$O$};
  \fill (m) circle (2.0pt) node[above right] {$m$};
  \draw[<->] (-2.35,-0.36) -- (2.35,-0.36) node[midway,below] {$L$};
  \draw[->,thick] (-1.62,-0.88) .. controls (-2.55,-0.65) and (-2.60,0.18) .. (-2.18,0.40);
  \draw[->,thick] (1.62,0.88) .. controls (2.55,0.65) and (2.60,-0.18) .. (2.18,-0.40);
\end{tikzpicture}

\caption{第37届复赛第7题双黑洞绕质心旋近的TikZ重绘。}\label{fig:cpho37q07}
\end{figure}

（1）设辐射功率与 $G$ 以及 $I$、$\omega$、$c$ 有关，推导量纲表达式。（2）初始间距为 $L_0$，求绕质心转动一周的时间。（3）求间距减为 $L_0/2$ 时已转过的角度。

\subsection{参考解答整理}
\paragraph{（1）量纲分析。} 令
\[
P=\alpha GI^\lambda\omega^\xi c^\sigma.
\]
由 $[P]=\mathrm{ML^2T^{-3}}$ 得
\[
\lambda=2,\qquad \xi=6,\qquad \sigma=-5,
\]
故本题近似模型为
\[
\boxed{P=\frac{\alpha GI^2\omega^6}{c^5}.}
\]
该式只用于题设的量纲分析和牛顿近似，不是完整的广义相对论辐射公式。

\paragraph{（2）转动一周所需时间。} 近似圆周运动满足
\[
\omega^2=\frac{G(M+m)}{L^3},\qquad I=\frac{Mm}{M+m}L^2,
\]
\[
E=-\frac{GMm}{2L}=-\frac{G^{2/3}Mm}{2(M+m)^{1/3}}\omega^{2/3}.
\]
由 $P=-\mathrm dE/\mathrm dt$ 得
\[
\omega^{-11/3}\frac{\mathrm d\omega}{\mathrm dt}=\frac{3\alpha G^{5/3}Mm}{c^5(M+m)^{1/3}},
\]
\[
\omega(t)=\left[\omega_0^{-8/3}-\frac{8\alpha G^{5/3}Mm}{c^5(M+m)^{1/3}}t\right]^{-3/8}.
\]
令转角增加 $2\pi$，得到
\[
\boxed{t_{360^\circ}=\frac{c^5L_0^4}{8\alpha G^3Mm(M+m)}
\left\{1-\left[1-\frac{10\pi\alpha G^{5/2}Mm\sqrt{M+m}}{c^5L_0^{5/2}}\right]^{8/5}\right\}.}
\]

\paragraph{（3）间距减半时的转角。} 当 $L=L_0/2$，
\[
t=\frac{15c^5L_0^4}{128\alpha G^3Mm(M+m)},
\]
从初态起的转角为
\[
\boxed{\Delta\theta=\frac{(1-\sqrt2/8)c^5L_0^{5/2}}{5\alpha G^{5/2}Mm\sqrt{M+m}}\frac{180^\circ}{\pi}.}
\]
\subsection{转写依据}
题干、图7a和公式均依据本题题面转写稿、答案公式核对表和图件核对入口。图7a以TikZ重绘，不使用原题截图。本题第（1）问10分，第（2）问24分，第（3）问6分。
